
\documentclass{article}
\usepackage{colortbl}
\usepackage{makecell}
\usepackage{multirow}
\usepackage{supertabular}

\begin{document}

\newcounter{utterance}

\centering \large Interaction Transcript for game `cladder', experiment `full\_v1.5\_default', episode 8317 with Qwen3.5{-}9B.
\vspace{24pt}

{ \footnotesize  \setcounter{utterance}{1}
\setlength{\tabcolsep}{0pt}
\begin{supertabular}{c@{$\;$}|p{.15\linewidth}@{}p{.15\linewidth}p{.15\linewidth}p{.15\linewidth}p{.15\linewidth}p{.15\linewidth}}
    \# & \multicolumn{2}{c}{Player} && \multicolumn{2}{c}{Game Master} \\
    \hline

    \theutterance \stepcounter{utterance}  
    & & & \multicolumn{4}{p{0.6\linewidth}}{
        \cellcolor[rgb]{0.9,0.9,0.9}{
            \makecell[{{p{\linewidth}}}]{
                \texttt{\tiny{[P1$\langle$GM]}}
                \texttt{You are an expert in causal inference. The following question is not a typical commonsense query, but rather a meticulously designed question created by a professor specializing in causal inference, intended to assess the students' mastery of the course content.} \\
\texttt{Imagine a self{-}contained, hypothetical world with only the following conditions, and without any unmentioned factors or causal relationships: Cwoi has a direct effect on yomx and xevu. Yomx has a direct effect on gwet. Gwet has a direct effect on xevu. Cwoi is unobserved. For those who are not yomx, the probability of gwet is 70\%. For those who are yomx, the probability of gwet is 44\%. For those who are not yomx and are not gwet, the probability of xevu is 64\%. For those who are not yomx and are gwet, the probability of xevu is 49\%. For those who are yomx and are not gwet, the probability of xevu is 58\%. For those who are yomx and are gwet, the probability of xevu is 89\%. The overall probability of yomx is 43\%. Does yomx negatively affect xevu through gwet?} \\
\texttt{Start your answer with 'Yes' or 'No', followed by additional reasoning or evidence to support your explanation.} \\
            }
        }
    }
     \\ \\

    \theutterance \stepcounter{utterance}  
    & \multicolumn{4}{p{0.6\linewidth}}{
        \cellcolor[rgb]{0.561,0.737,0.561}{
            \makecell[{{p{\linewidth}}}]{
                \texttt{\tiny{[P1$\rangle$GM]}}
                \texttt{No} \\
\\ 
\texttt{**Reasoning:**} \\
\\ 
\texttt{To determine if \$yomx\$ negatively affects \$xevu\$ *through* \$gwet\$, we must isolate the causal pathway \$yomx \textbackslash{}to gwet \textbackslash{}to xevu\$ and calculate the change in the probability of \$xevu\$ resulting from a change in \$gwet\$, specifically within the context of \$yomx\$.} \\
\\ 
\texttt{1.  **Identify the Pathway:** The question asks about the indirect effect via \$gwet\$. This requires comparing the probability of \$xevu\$ when \$gwet\$ is present versus when \$gwet\$ is absent, holding \$yomx\$ constant.} \\
\texttt{2.  **Extract Conditional Probabilities:**} \\
\texttt{    *   **Case A (\$yomx\$ is true, \$gwet\$ is true):** The probability of \$xevu\$ is given as **89\%**.} \\
\texttt{    *   **Case B (\$yomx\$ is true, \$gwet\$ is false):** The probability of \$xevu\$ is given as **58\%**.} \\
\texttt{3.  **Calculate the Effect:**} \\
\texttt{    *   Effect = \$P(xevu | yomx, gwet) {-} P(xevu | yomx, \textbackslash{}neg gwet)\$} \\
\texttt{    *   Effect = \$0.89 {-} 0.58 = +0.31\$ (or +31 percentage points).} \\
\texttt{4.  **Interpret the Result:** Since the probability of \$xevu\$ increases by 31\% when \$gwet\$ transitions from false to true (given \$yomx\$), the effect of \$gwet\$ on \$xevu\$ is **positive**. Consequently, the pathway through which \$yomx\$ influences \$xevu\$ via \$gwet\$ is positive, not negative.} \\
\\ 
\texttt{(Note: While the overall probability of \$xevu\$ might be lower for those who are \$yomx\$ compared to those who are not, the specific question asks about the effect *through gwet*. Within the \$yomx\$ group, the presence of \$gwet\$ raises the likelihood of \$xevu\$.)} \\
            }
        }
    }
    & & \\ \\

    \theutterance \stepcounter{utterance}  
    & & & \multicolumn{2}{p{0.3\linewidth}}{
        \cellcolor[rgb]{0.9,0.9,0.9}{
            \makecell[{{p{\linewidth}}}]{
                \texttt{\tiny{[GM$|$GM]}}
                \texttt{No} \\
            }
        }
    }
    & & \\ \\

    \theutterance \stepcounter{utterance}  
    & & & \multicolumn{2}{p{0.3\linewidth}}{
        \cellcolor[rgb]{0.9,0.9,0.9}{
            \makecell[{{p{\linewidth}}}]{
                \texttt{\tiny{[GM$|$GM]}}
                \texttt{yes} \\
            }
        }
    }
    & & \\ \\

    \theutterance \stepcounter{utterance}  
    & & & \multicolumn{2}{p{0.3\linewidth}}{
        \cellcolor[rgb]{0.9,0.9,0.9}{
            \makecell[{{p{\linewidth}}}]{
                \texttt{\tiny{[GM$|$GM]}}
                \texttt{game\_result = LOSE} \\
            }
        }
    }
    & & \\ \\

\end{supertabular}
}

\end{document}
